% ===================================================================
%  도표 제 14 호 — 거리. 장사하는 사람들.
%
%  Traduction, faite en 2026, en coreen d'aujourd'hui, sur l'ido de
%  texto/io. Regles de la colonne : voir 10-dopyo-01.tex. Vingt-deux
%  alineas, quatre-vingt-dix-neuf renvois — le tableau le plus long
%  du livret apres le 5.
%
%  ET SEIZE PAIRES SEULEMENT, contre trente-trois au tableau 13 qui
%  compte six renvois de MOINS. Le 13 attache, le 14 enumere
%  — le marchand de parapluies, le teinturier, l'armurier, les
%  fiacres, les voitures de maitre — et deux renvois joints par « e »
%  ne sont pas en rapport de dependance.
%
%  ET ON NE MELERA PAS LES DEUX UNITES. Le marathe avait releve neuf
%  inversions contre trois a ces deux tableaux, le telougou dix-neuf
%  contre neuf : ce sont des inversions REELLEMENT faites. Seize et
%  trente-trois sont des paires LISTEES par parigi.py, et ce compte-la
%  ne mesure pas la colonne coreenne — il mesure l'ido, puisque
%  parigi.py ne lit que l'ido et sort la meme liste pour toutes les
%  colonnes. Le rapport de deux pour un vaut dans les deux unites, et
%  c'est bien la meme cause ; mais la colonne coreenne n'apporte pas
%  ici un troisieme releve independant, elle confirme dans une
%  troisieme famille que la depense est dans la syntaxe de 1926.
%
%  LES METIERS DE CASTE, TROISIEME FOIS DANS CETTE COLONNE. Le
%  cordonnier du tableau 7 avait deja pose la question ; ici
%  reviennent le coiffeur (23) et le teinturier (53). Le coreen a
%  백정 pour le boucher et 갖바치 pour le cordonnier — deux des
%  metiers 천민, encore employes comme insultes. Le coiffeur, lui,
%  n'a pas ce probleme : 이발사 est le mot ordinaire des enseignes.
%  Le teinturier s'ecrit 염색하는 사람, la fonction. Meme regle,
%  meme resultat, et c'est la troisieme fois qu'elle tient sans qu'on
%  ait a la reinventer.
%
%  L'IMPRIMERIE ET L'HORLOGERIE SONT ENTIERES EN COREEN : 인쇄소
%  l'imprimerie, 식자공 les typographes, 석판 인쇄 la lithographie,
%  금은 세공 l'orfevrerie, 재봉틀 la machine a coudre, 굴뚝 청소부 le
%  ramoneur, 소매상 le mercier. Le cycle, l'automobile et la
%  photographie ne le sont pas : 사진사, 타자수, 자전거 상인 —
%  la coupure passe encore a l'annee et non a la langue.
% ===================================================================

%%K t14-apar-1 apar
\VUsaut{4.0mm}
\VUtitre{15.0pt}{60}{도표 제 14 호}
\VUsaut{3.0mm}

%%K t14-tit-1 sub
\VUpk{11.4pt}{40}{거리. --- 장사하는 사람들.}
\VUsaut{3.0mm}

%%K t14-01-1 p
1. --- 시골에 사는 내 벗 요안네스가 우리 집에서 사흘을 지내러 왔다. \VUgras{그때까지}
그는 큰 도시를 볼 기회가 없었다. 그래서 우리는 날마다 거닐며 지내고, 나는 그가 모르는
것을 설명해 준다.

%%K t14-02-1 p
2. --- 먼저 \VUgras{천문대} \textsuperscript{(1)}를 보러 갔다. 그가 아주 흥미로워했다.
돌아오는 길에 \VUgras{거리} \textsuperscript{(3)}를 지났다 — 거기에 \VUgras{터키탕}
\textsuperscript{(2)}이 있다 —, 거기에서 \VUgras{장례 행렬} \textsuperscript{(4)}을 만났다.
\VUgras{상여 행렬} 앞에는 \VUgras{성직자들}이 걸었고, \VUgras{영구차} 앞에 섰다. 그
차에 \VUgras{널}을 실었다. 많은 벗이 \VUgras{상을 당한} 집안과 함께 걸었다.

%%K t14-02-2 p
행렬이 지나간 뒤 우리는 큰 \VUgras{맥줏집} \textsuperscript{(5)} 앞에서 길을 건넜다. 그
집 옥상에서는 많은 \VUgras{손님}이 \VUgras{맥주}를 마시고 있었다 — 유리를 끼운
\VUgras{차양} \textsuperscript{(6)} 아래에서.

%%K t14-03-1 p
3. --- \VUgras{술 파는 사람} \textsuperscript{(7)}의 가게와 \VUgras{악기 상인}
\textsuperscript{(10)}의 가게 사이에 걸린 \VUgras{방패 표지}를 보고 요안네스가 몹시
놀랐다. 그것이 무슨 뜻이냐고 물었다. 나는 그것이 영국 \VUgras{영사관} \textsuperscript{(8)}
을 알린다고 대답하고, 발코니에 있는 \VUgras{영사} \textsuperscript{(9)}를 그에게 가리켜
주었다.

%%K t14-tit-2 sub
\VUpk{10.2pt}{40}{가게 구경.}
\VUsaut{3.0mm}

%%K t14-04-1 p
4. --- 물론 우리는 \VUgras{악기}와 \VUgras{하모늄}과 \VUgras{오르간}과
\VUgras{피아노} 진열 앞에서 멈추었다. \VUgras{악보}와 피아노 \VUgras{곡집}의 그림 표지를
보았고, 금박을 입힌 나무로 만든 \VUgras{수금}을 감탄했다 — 그것을 간판으로 쓴다.

%%K t14-05-1 p
5. --- \VUgras{길 쓰는 사람} \textsuperscript{(11)}들과 \VUgras{청소 수레} \textsuperscript{(12)}
가 일으킨 먼지구름 때문에 우리는 서둘러 지나가야 했다 — \VUgras{가구 상인} \textsuperscript{(13)}
의 고운 \VUgras{세간} 앞을, 그리고 \VUgras{양철장이} \textsuperscript{(14)}의 \VUgras{화덕}과
\VUgras{등} 진열 앞을.

%%K t14-06-1 p
6. --- 게다가 여기에서는 인도를 벗어나야 했다. 인도에 꾸러미가 잔뜩 놓여 있었기
때문이다. 그것을 \VUgras{짐수레} \textsuperscript{(16)}에서 내려, 방금 세를 얻은
\VUgras{가게} \textsuperscript{(15)}에 들여놓으려는 것이었다.

%%K t14-06-2 p
그 바람에 \VUgras{수레 만드는 이} \textsuperscript{(17)}의 고운 수레들은 보지 못했다.
그러나 \VUgras{짐 꾸리는 이} \textsuperscript{(19)}를 보려고는 멈추어 섰다. 그는 제
이웃에게 줄 궤를 만들고 있었는데, 그 이웃은 \VUgras{자전거와 자동차 상인} \textsuperscript{(18)}
이다. 그 다음, 좀 늦은 터라 우리는 집으로 돌아왔다.

%%K t14-tit-3 sub
\VUpk{10.2pt}{40}{도시 거닐기.}
\VUsaut{3.0mm}

%%K t14-07-1 p
7. --- 이튿날 어머니가 요안네스의 사진을 찍으라고 하시고, 나에게 \VUgras{사진사}
\textsuperscript{(20)}에게 함께 가라고 이르셨다. 점심 뒤에 우리는 그 예술가의
\VUgras{작업실}로 갔다. 거기에서 꽤 오래 기다려야 했다. 심심풀이로 벽에 걸린
\VUgras{초상} \textsuperscript{(22)}들을 보았다.

%%K t14-07-2 p
우리 차례가 되었을 때 일은 오래 걸리지 않았다. 내 벗이 \VUgras{사진기} \textsuperscript{(21)}
앞에 서기를 마치자마자 우리는 내려왔다.

%%K t14-08-1 p
8. --- 이 층에서 \VUgras{이발사} \textsuperscript{(23)}의 문이 열려 있어, 그의 조수가
손님의 머리를 자르는 것이 보였다.

%%K t14-08-2 p
거리로 나와 우리는 \VUgras{담배 가게} \textsuperscript{(24)} 앞을 지났다. 붉은
\VUgras{초롱} \textsuperscript{(25)}과, \VUgras{간판} \textsuperscript{(26)}으로 쓰는 굵은
\VUgras{담뱃대}가 요안네스를 어찌나 즐겁게 했는지, 그는 \VUgras{코담배를 맡으며}
\textsuperscript{(27)} 이야기하던 늙은 어른 둘과 부딪치고 말았다.

%%K t14-tit-4 sub
\VUpk{10.2pt}{40}{과자 가게.}
\VUsaut{3.0mm}

%%K t14-09-1 p
9. --- \VUgras{환전상} \textsuperscript{(29)}의 \VUgras{진열창}에 늘어놓은 온 나라의
\VUgras{지폐}와 \VUgras{동전}이 잠시 우리 눈길을 끌었다. 그러나 새참 때가 되어 우리는
더 머물지 않고 \VUgras{과자 장수} \textsuperscript{(31)}의 가게로 들어갔다.

%%K t14-10-1 p
10. --- 가게에 있는 온갖 \VUgras{맛난 것}을 보고 — \VUgras{과자}, \VUgras{꿀 과자},
\VUgras{비스킷}, 온갖 \VUgras{사탕} — 우리는 쉽게 고르지 못했다. 마침내 요안네스는
\VUgras{크림 과자}를 골랐고, 나는 작은 \VUgras{버찌 타르트}만 집었다. 단것은
\VUgras{내 이를 시게 하고}, \VUgras{치과 의사} \textsuperscript{(30)}에게 자주 가고 싶은
마음은 조금도 없기 때문이다.

%%K t14-tit-5 sub
\VUpk{10.2pt}{40}{기계로 쓰기.}
\VUsaut{3.0mm}

%%K t14-11-1 p
11. --- 이야기만 듣고 본 적은 없는 \VUgras{타자기}를 내 벗에게 보여 주려고, 우리는
\VUgras{사무실} \textsuperscript{(32)}로 들어갔다 — \VUgras{내 사촌 형} \textsuperscript{(34)}
의 사무실이다. 그가 제 편지를 \VUgras{비서} \textsuperscript{(35)}에게 부르는 것을 보았다.
그 비서는 \VUgras{속기하는 사람}이다. 편지가 끝나자마자 \VUgras{타자수} \textsuperscript{(33)}
가 그것을 제 기계로 옮겨 쳤다. 친척의 일을 방해하고 싶지 않아 우리는 잠깐만 머물렀다.

%%K t14-12-1 p
12. --- 사촌 형의 사무실 아래에는 큰 \VUgras{시계와 보석 상인} \textsuperscript{(37)}이
산다. 그는 금은 세공사이기도 하다. 그의 훌륭한 가게에는 많은 \VUgras{괘종시계},
\VUgras{추시계}, \VUgras{회중시계}, \VUgras{작은 사슬}, 값진 \VUgras{보석}이 있다 —
이를테면 \VUgras{진주 목걸이}, 금 \VUgras{팔찌}, \VUgras{귀고리}, \VUgras{보석}과
\VUgras{금강석}을 박은 \VUgras{반지}다. 진열창 하나는 \VUgras{금붙이}만을 위해 따로
두었고, 거기에 은 \VUgras{식기}가 크게 모여 있다.

%%K t14-13-1 p
13. --- 거리 모퉁이를 돌면서 나는 집 \VUgras{번지} 곁에 붙인 \VUgras{문패} \textsuperscript{(36)}
가 바뀐 것을 알아차렸다. 내 말을 크게 하려던 참에 군악의 즐거운 소리가 들렸다. 우리는
연대 쪽으로 달려갔다 — 그 연대가 \VUgras{모자 상인} \textsuperscript{(39)}과
\VUgras{장갑 상인} \textsuperscript{(40)}의 가게 앞을 지나갈 것이고, \VUgras{안경 상인}
\textsuperscript{(38)}의 진열창 앞에 서 있어도 행진이 그만큼 잘 보였으리라는 것은
생각하지 않고. 그 진열창에는 \VUgras{코안경}, \VUgras{귀걸이 안경},
\VUgras{작은 망원경}이 그득하다.

%%K t14-tit-6 sub
\VUpk{10.2pt}{40}{행진.}
\VUsaut{3.0mm}

%%K t14-14-1 p
14. --- \VUgras{연대} \textsuperscript{(41)} 앞에서 \VUgras{북 대장} \textsuperscript{(42)}이
걸었다 — 그 뒤로 \VUgras{북 치는 이} \textsuperscript{(43)}들과 \VUgras{나팔 부는 이}
\textsuperscript{(44)}들과 \VUgras{악대} 전체가 따랐다. \VUgras{장군} \textsuperscript{(45)}은
제 부관과 함께 광장에 있었다. 그는 \VUgras{물줄기} \textsuperscript{(49)} 가까이에 서 있었다
— \VUgras{분수} \textsuperscript{(48)}의 물줄기다 —, \VUgras{행진}을 보려는 것이다.

%%K t14-14-2 p
\VUgras{참모 장교} \textsuperscript{(46)}들과 \VUgras{대령}과 \VUgras{중령}이 긴 칼로 그에게
경례했다. \VUgras{소령}들은 제 \VUgras{대대} \textsuperscript{(47)} 앞에 있었고, 대위마다
제 \VUgras{중대}를 이끌었다. 그동안 \VUgras{중위}들과 \VUgras{소위}들과
\VUgras{부사관}들은 제 부하 곁 늘 있던 자리에 있었다.

%%K t14-15-1 p
15. --- 많은 사람이 \VUgras{네거리} \textsuperscript{(50)}에 모였다. 장사하는 사람들,
\VUgras{우산 상인} \textsuperscript{(51)}, \VUgras{염색하는 사람} \textsuperscript{(53)},
\VUgras{무기 상인} \textsuperscript{(54)}이 \VUgras{병사}들을 보려고 나왔다. 모든 탈것이 —
\VUgras{삯마차} \textsuperscript{(52)}, \VUgras{주인 마차} \textsuperscript{(57)},
\VUgras{물 뿌리는 통 수레} \textsuperscript{(56)}까지 — 인도를 따라 비켜섰다.

%%K t14-tit-7 sub
\VUpk{10.2pt}{40}{거리의 장면.}
\VUsaut{3.0mm}

%%K t14-15-2 p
\VUgras{돌아다니는 장사꾼} \textsuperscript{(55)} 하나가 손에 제 \VUgras{견본 상자}를 들고
빠르게 길을 건넜다. \VUgras{윗층 자리} \textsuperscript{(59)}에 앉은 사람들이 —
\VUgras{합승 마차} \textsuperscript{(58)}의 윗층이다 — 구경거리를 보려고 몸을 돌렸다. 딱한
\VUgras{떠돌이 노래꾼} \textsuperscript{(60)} 하나가 제 \VUgras{손풍금} \textsuperscript{(61)}
과 함께 제 \VUgras{노래}를 멈추어야 했다 — 북소리가 그의 목소리를 덮었기 때문이다.

%%K t14-16-1 p
16. --- 연대가 \VUgras{옷감 가게} \textsuperscript{(64)} 앞에 이르자, \VUgras{바느질하는
여자} \textsuperscript{(63)}들과 \VUgras{재봉사} \textsuperscript{(62)}들이 제
\VUgras{재봉틀}과 제 일을 놓고 창으로 내다보았다. \VUgras{그 상인} \textsuperscript{(65)}은
— 방금 새 \VUgras{옷} \textsuperscript{(66)}을 제 \VUgras{진열대}에 걸어 놓은 사람이다 —
인도에 내놓은 \VUgras{휘장} \textsuperscript{(67)}과 \VUgras{삼베}를 안으로 들여놓아야
했다. 그것들은 \VUgras{하수구 뚜껑} \textsuperscript{(68)} 가까이에 있었다. 사람들이 그것을
넘어뜨릴까 걱정한 것이다. 늙은 \VUgras{봇짐 장수} \textsuperscript{(69)}까지도 — 그에게서
어떤 문지기 아낙이 양말을 흥정하고 있었다 — 제 흥정을 마치려고 복도로 들어가야 했다.

%%K t14-16-2 p
우리는 연대를 병영까지 따라갔다. 그날에 아주 흡족해서 저녁 시간에 맞추어 집으로 돌아왔다.

%%K t14-tit-8 sub
\VUpk{10.2pt}{40}{여러 장사와 생업.}
\VUsaut{3.0mm}

%%K t14-17-1 p
17. --- 이튿날 아버지가 그곳 신문의 인쇄소를 보러 가자고 하셨다. 우리는 신이 나서
그러자고 했다.

%%K t14-17-2 p
\VUgras{인쇄인}은 \VUgras{책 상인} \textsuperscript{(70)} \VUgras{(= 책 파는 사람)}이기도
하고, \VUgras{출판인}, \VUgras{종이 상인}, \VUgras{제책하는 사람}이기도 하다. 그는
이 층에 신문 \VUgras{편집실} \textsuperscript{(71)}을 두었고, \VUgras{조각실}도 두었다 —
거기에서 \VUgras{석판 인쇄공} \textsuperscript{(72)}들과 \VUgras{구리 조각공} \textsuperscript{(73)}
들이 일한다 —, 끝으로 \VUgras{식자실}도 두었다 — 거기에 \VUgras{인쇄 직공} \textsuperscript{(74)}
들이 있다. 참된 \VUgras{인쇄소} \textsuperscript{(75)}와 \VUgras{인쇄기}와 여러 기계는
아래층에 있다.

%%K t14-tit-9 sub
\VUpk{10.2pt}{40}{요안네스가 묻는 것이 아주 많다.}
\VUsaut{3.0mm}

%%K t14-18-1 p
18. --- 인쇄소에서 나오면서 요안네스는 몹시 놀라 기둥 위에 놓인 작은 유리 상자 앞에
멈추어 섰다 — \VUgras{소매 가게} \textsuperscript{(77)} 맞은편이다 —, 그러고는 그것이
무엇에 쓰이느냐고 물었다. 아버지가 그것은 \VUgras{화재 경보기} \textsuperscript{(76)}
라고 설명해 주셨다. 사실 도시의 모든 것이 내 벗에게는 놀라움이었다 — 수수한
\VUgras{돌 우물} \textsuperscript{(78)}에서부터, 가벼운 \VUgras{세발 짐수레} \textsuperscript{(80)}
에서부터 — 그것을 제복 입은 \VUgras{사환} \textsuperscript{(79)}이 몬다 —
\VUgras{우편 수레} \textsuperscript{(81)}까지.

%%K t14-19-1 p
19. --- 아버지가 잠시 우리를 두고 \VUgras{세무 관리} \textsuperscript{(83)}와 이야기하러
가셨다 — 그는 \VUgras{변호사} \textsuperscript{(82)}와 \VUgras{공증인} \textsuperscript{(84)}
와 이야기하려고 멈추어 서 있었다 —, 그동안 우리는 거리의 오감을 눈으로 좇으며 심심풀이를
했다. 우리 앞으로 아주 여러 사람이 지나갔다. \VUgras{과자 가게 사환} \textsuperscript{(85)}
하나가 바구니에 훌륭한 \VUgras{고기 파이}를 이고 \VUgras{헌옷 장수} \textsuperscript{(86)}
뒤에 왔다. \VUgras{가로등 켜는 이} \textsuperscript{(87)}가 \VUgras{가스 일꾼} \textsuperscript{(88)}
과 이야기하고 있었다. \VUgras{수녀} \textsuperscript{(89)} 하나가 고아 소녀의 손을 잡고 제
수도원으로 돌아가고 있었다.

%%K t14-tit-10 sub
\VUpk{10.2pt}{40}{집으로 돌아오면서.}
\VUsaut{3.0mm}

%%K t14-20-1 p
20. --- 아버지가 우리와 다시 만나신 순간, 요안네스가 크게 웃었다. 온통 검댕을 뒤집어쓴
\VUgras{굴뚝 청소부} \textsuperscript{(91)} 둘의 더러운 얼굴과 이상한 누더기를 보고서다.

%%K t14-21-1 p
21. --- 나는 어머니께 드릴 화초를 \VUgras{꽃 파는 여자} \textsuperscript{(92)}에게서
사고 싶었다. 그러나 그것은 들고 가기에 너무 무거우니 \VUgras{오렌지}를 사는 편이 낫겠다고
아버지가 일러 주셨다. 우리는 \VUgras{파는 여자} \textsuperscript{(94)}에게 다가갔다.
\VUgras{가정 교사} \textsuperscript{(93)}가 — 그 여자는 제 어린 제자에게 줄 것을 고르고
있었다 — 사기를 마치자 우리도 사서 받았다.

%%K t14-22-1 p
22. --- 집으로 돌아오면서 우리는 아는 이를 여럿 만났다. 우리 집안의 늙은 벗 부인은 —
그분은 제 \VUgras{말벗 부인} \textsuperscript{(95)}의 팔에 손을 얹고 있었다 —, 다리와 길의
\VUgras{기술자} \textsuperscript{(96)}, 그리고 내 사촌 누이 하나가 제 \VUgras{아이 보는 여자}
\textsuperscript{(98)}와 함께였다.

%%K t14-22-2 p
그 다음에 우리는 어머니의 \VUgras{모자 만드는 여자} \textsuperscript{(97)}를 만났고, 도시에
낯선 \VUgras{수도사} \textsuperscript{(99)} 둘도 만났다 — 그들은 정거장으로 가는 길을 묻고자
아버지께 다가왔다.

\VUsaut{6.0mm}
\VUfilet{22mm}
